\section{文本编辑}
\label{sec:nlp-generation-editing}

前几节的输出几乎都要从头写出；\term{文本编辑}（text editing）则从一段已经存在的文字出发，
输出与输入的大部分词元相同。改写调整表达，风格转换改变指定属性，纠错修复错误，
简化为特定读者降低理解负担。四类任务共有四个挑战：保留不该改的内容，
多数位置的正确动作是原样复制；找到该改的位置，修改往往稀疏而分散；
缺少同内容不同表达的平行句对，风格转换几乎没有；区分漏改、正确修改、过改和意义改变，
把“须提交”改成“不必提交”只动了几个字，意义变化却大于一次保持原意的整句重写。

方法谱系沿两条线展开。一条是输出的表示方式：规则与统计替换只在少数位置提出候选；
整句序列生成用第~\ref{sec:nlp-generation-encoder-decoder}小节的编码器—解码器重写全句，
借第~\ref{sec:nlp-generation-copy-pointer}小节的复制保留原词；编辑标签方法只为每个源词元
预测一个动作；大语言模型以提示完成编辑。另一条是监督来源：有平行句对时直接训练，
没有时用回译、属性语料、合成错误或挖掘复述构造训练信号。四个小节按修改理由划分任务，
在每个任务内沿这两条线比较方法，节末表~\ref{tab:nlp-generation-editing-methods}汇总对照。

\paragraph{从源与输出的对齐恢复修改}
上面两条线都要回答同一个问题：一对句子之间究竟改了什么。编辑标签方法要据此把平行句对
转成逐词元的训练标签，整句生成要据此从输出中找出实际修改，评价也要在修改集合上计数。
\term{编辑距离}（edit distance）是把源序列$x_{1:n}$变成目标序列$y_{1:m}$所需的
最少插入、删除和替换次数；用动态规划求它并回溯路径，就得到两者的
\term{Levenshtein对齐}。设$d_{i,j}$为前缀$x_{1:i}$与$y_{1:j}$之间的编辑距离，
以$d_{i,0}=i$、$d_{0,j}=j$初始化，按行递推
\begin{equation}
  d_{i,j}=\min\bigl\{\,d_{i-1,j}+1,\;\; d_{i,j-1}+1,\;\;
  d_{i-1,j-1}+\mathbb{I}[x_i\neq y_j]\,\bigr\}.
  \label{eq:nlp-generation-edit-distance}
\end{equation}
三项依次对应删除$x_i$、插入$y_j$、替换或保留当前位置。递推只依赖左、上和左上三个格子，
时间与空间都是$O(nm)$；记录每格取到最小值的来源，从$(n,m)$回溯到$(0,0)$，
就得到操作序列及每个操作的源坐标。几个来源并列最小时，按预先固定的优先次序取一个。

以$x=$“go／to／meeting”、$y=$“goes／to／the／meeting”为例，
表~\ref{tab:nlp-generation-levenshtein}给出全部$d_{i,j}$。右下角$d_{3,4}=2$；
回溯时，$(3,4)$来自左上方的匹配meeting，$(2,3)$来自左方，即插入the，
$(2,2)$来自左上方的匹配to，$(1,1)$来自左上方的替换go为goes。
两次编辑恰好是一次词形修正和一次冠词插入。

\begin{table}[htbp]
  \centering
  \small
  \begin{tabular}{@{}lccccc@{}}
    \toprule
    & $\varnothing$ & goes & to & the & meeting\\
    \midrule
    $\varnothing$ & 0 & 1 & 2 & 3 & 4\\
    go & 1 & \textbf{1} & 2 & 3 & 4\\
    to & 2 & 2 & \textbf{1} & \textbf{2} & 3\\
    meeting & 3 & 3 & 2 & 2 & \textbf{2}\\
    \bottomrule
  \end{tabular}
  \caption{“go to meeting”到“goes to the meeting”的编辑距离表。行是源前缀，列是目标前缀，
  粗体格是回溯路径：从右下角经匹配、插入、匹配回到替换，共两次编辑。}
  \label{tab:nlp-generation-levenshtein}
\end{table}

若去掉替换操作、只允许插入和删除，编辑距离的回溯路径就是最长公共子序列对齐，
后文LaserTagger构造标签用的正是这一特例；GECToR用Levenshtein对齐的一个变体，
第~\ref{sec:nlp-generation-edit-metrics}小节的编辑指标也在对齐恢复的编辑集合上计数。
对齐只说明改了哪里，不说明改得对不对：删除“不得参加”中的“不”同样只需一次编辑。

\subsection{改写与复述}
\label{sec:nlp-generation-paraphrase}

\term{改写}（paraphrasing）要求意义不变而表达变化。以通知为例，把“须在4月16日18:00前提交报名信息”
改成“必须在4月16日18:00前提交报名信息”可以接受；改成“可在4月16日提交信息”
同时丢掉了义务、精确时点和“前”。改写有两个核心困难：训练需要大量意义相同的句对，
即\term{复述}（paraphrase），而人工写复述代价高；句子高度重合时，整句生成仍要逐词重写，
既浪费计算，也给补造内容留下机会。回译批量构造复述对，回应前一点；
局部编辑标注只预测改动，回应后一点。

\subsubsection{回译复述}
\label{sec:nlp-generation-backtranslation-paraphrase}

两句话若能译成同一句外语，它们多半表达相同的意思。
Bannard与Callison-Burch据此从双语平行语料中抽取短语复述：
设$e_1,e_2$为同一语言的两个短语，$f$为与它们对齐的外语短语，
复述概率按枢轴语言上的全概率展开
\begin{equation}
  p(e_2\mid e_1)=\sum_{f}p(e_2\mid f)\,p(f\mid e_1),
  \label{eq:nlp-generation-pivot-paraphrase}
\end{equation}
两个条件概率都由短语表的对齐计数估计\scite{bannard2005paraphrasing}
式~\eqref{eq:nlp-generation-pivot-paraphrase}把“意义相同”换成“共享译文”，
其中求和覆盖了一个短语的多种译法，单个错误对齐的影响因此被其他译法稀释。

神经翻译把这一思想推广到整句：把句子$x$译成枢轴语言$z$，再把$z$译回原语言，
得到复述$\tilde x$。Mallinson等用多条枢轴译文及多个枢轴语言对回译分布加权平均，
近似式~\eqref{eq:nlp-generation-pivot-paraphrase}的求和\scite{mallinson2017paraphrasing}
Wieting与Gimpel走得更远：他们取大规模捷克语—英语平行语料，把捷克语一侧机器翻译成英语，
与原有英语句配成五千万量级的复述对ParaNMT，再按句向量相似度和词重合度过滤，
用来训练复述句向量\scite{wieting2018paranmt}
这些做法只需现成的翻译系统或双语语料，不需要人工复述。

回译的多样性取决于译回时怎样选择候选。束搜索倾向于给出最常见的译法，
结果常与原句几乎相同；从模型分布中采样或在束内引入随机性，
表达更多样，也更容易偏离原意。因此要过滤：相似度太低的句对可能已经改变意义，
太高的又几乎没有改写。设$\operatorname{sim}$为句向量余弦相似度，可只保留
$\tau_{\mathrm{low}}\leq\operatorname{sim}(x,\tilde x)\leq\tau_{\mathrm{high}}$且编辑距离大于零的句对。
整句相似度对少数关键词不敏感：通知经英语回译成“来宾可在4月16日提交信息”时，
相似度可能仍然不低，却丢掉了人群、义务和时点。日期、实体和否定词这类保护片段
应在过滤阶段逐项比对。

回译复述的长处是规模大、成本低、领域可以跟随单语材料；
代价是意义漂移只能事后过滤，而且复述的多样性受限于翻译系统本身的表达习惯。

\subsubsection{局部编辑标注：LaserTagger}
\label{sec:nlp-generation-lasertagger}

有了复述对，可以用编码器—解码器整句生成输出。但在句子融合、拆句和局部改写里，
输出往往只在几个位置与输入不同，整句生成需要逐词重写那些本该照抄的部分。
LaserTagger把输出表示为对源词元的一串编辑动作，只需预测这些动作，
再由确定的程序重建输出\scite{malmi2019lasertagger}

\paragraph{标签与短语词表}
LaserTagger为每个源词元$x_i$分配一个标签$t_i=(b_i,p_i)$：基础动作
$b_i\in\{\code{KEEP},\code{DELETE}\}$决定是否保留$x_i$，可选的插入短语$p_i$
放在$x_i$之前，取值于\term{短语词表}（phrase vocabulary）$V$或为空。
标签集的大小是$2(|V|+1)$。替换由“删除原词并在它之前插入新短语”表示。
以句子融合为例，源文本
“Turing／was／born／in／1912／.／Turing／died／in／1954／.”的目标是
“Turing was born in 1912 and died in 1954.”，只需给第一个句号打\code{DELETE}，
给第二个Turing打\code{DELETE}，再给died打$(\code{KEEP},\text{and})$，即在died之前插入and，
其余全部\code{KEEP}。
中文通知中“须”改为“必须”，同样是在“须”上打$(\code{DELETE},\text{必须})$。

短语词表决定了模型能表达哪些输出。把训练集每对句子用最长公共子序列对齐，
未被对齐的目标片段就是需要插入的短语。理想的$V$应在大小上限$k$内
使尽可能多的训练样本可以完整重建，这个选择问题是NP困难的；
LaserTagger按短语在训练集中被需要的次数排序，取前$k$个。
例如三对训练句分别需要插入$\{\text{and}\}$、$\{\text{and}\}$和$\{\text{, which}\}$，
取$k=1$时选and，覆盖三对中的两对；第三对无法用标签表示，训练时被丢弃。

\paragraph{标签转换与模型}
给定平行句对和$V$，下面的贪心程序把目标转成标签；转换失败的样本不参与训练。

\begin{algorithm}[LaserTagger的标签转换]
  \label{alg:nlp-generation-lasertagger-convert}
  \begin{algorithmic}[1]
    \Require 源词元$x_{1:n}$，目标词元$y_{1:m}$，短语词表$V$
    \Ensure 标签序列$t_{1:n}$，或“无法表示”
    \State $j\gets1$\Comment{目标中下一个待匹配的位置}
    \For{$i=1,\ldots,n$}
      \If{存在$j'\geq j$使$y_{j'}=x_i$，且$y_{j:j'-1}$为空或属于$V$}
        \State 取满足条件的最小$j'$；$t_i\gets(\code{KEEP},y_{j:j'-1})$；$j\gets j'+1$
      \Else
        \State $t_i\gets(\code{DELETE},\varnothing)$
      \EndIf
    \EndFor
    \If{$j\leq m$}\Comment{目标末尾仍有未匹配词元}
      \State 若$y_{j:m}\in V$，把它作为插入短语附到末尾的特殊位置；否则\Return 无法表示
    \EndIf
  \end{algorithmic}
\end{algorithm}

对上面的句子融合，算法~\ref{alg:nlp-generation-lasertagger-convert}处理第一个句号时，
目标中下一个句号在句末，两者之间的“and died in 1954”不在词表中，打\code{DELETE}；第二个Turing在目标剩余部分中不出现，也打\code{DELETE}；
died与目标匹配，其前的未匹配片段and在词表中，于是得到$(\code{KEEP},\text{and})$。
模型用BERT编码源句，得到$\bm H\in\mathbb{R}^{n\times768}$。
前馈版本LaserTagger\textsubscript{FF}对每个位置独立分类：
$p(t_i\mid x)=\operatorname{softmax}(\bm W\bm h_i+\bm b)$，$\bm W\in\mathbb{R}^{|T|\times768}$，
$|T|$为标签数；自回归版本LaserTagger\textsubscript{AR}在编码器上接一层Transformer解码器，
逐位置预测标签，并以已预测标签为条件。两者都最小化逐位置交叉熵
$-\sum_{i}\log p(t_i\mid x,t_{<i})$，前馈版本中条件$t_{<i}$不出现。
推断时把标签依次应用到源句即可得到输出。

\paragraph{关联、推广与边界}
LaserTagger是第~\ref{chap:nlp-text-analysis}章序列标注在生成任务中的用法：
输出空间从“任意词元序列”收缩为“每个源词元一个标签”，前馈版本可以并行预测全部位置，
推断比同规模的序列生成快得多。它的约束也来自这里：插入只能取词表内的短语，
大幅重排无法表达，训练集中需要罕见插入的样本只能丢弃。
Seq2Edits把标签推广为带跨度的编辑序列，每步输出“编辑类型、跨度终点、替换短语”三元组，
替换短语可从完整词表生成，从而能表达跨多个词的修改\scite{stahlberg2020seq2edits}
在源与目标高度重合、需要插入的短语种类有限时，编辑标签方法通常比整句生成更快，
也更少补造内容；修改幅度越大，整句生成越有优势。

\subsection{无平行数据的风格转换}
\label{sec:nlp-generation-style-transfer}

\term{风格转换}（style transfer）要改变句子的一个属性，例如情感、正式程度或礼貌程度，
同时保留其余内容。形式上，给定按属性分开的单语语料$X_1,\ldots,X_K$，
没有同内容不同属性的句对，要学习映射$x\mapsto\tilde x$，使$\tilde x$具有目标属性$a'$，
并与$x$讲同一件事。缺少平行数据是这项任务的核心挑战，各方法的差别在于用什么替代平行监督：
DRG假定属性集中在少数可删除的短语，用检索与重建代替；
隐表示方法试图把内容与属性分开，或改用回译直接构造伪平行对；
大语言模型借预训练中学到的风格知识，以提示直接转换。
评价也因此分三项：目标属性是否达成，通常用属性分类器判断；
内容是否保持，用与源句的相似度或人工判断；输出是否流畅。
三项相互牵制：原样复制源句的内容保持最好，属性却没有改变。

\subsubsection{Delete--Retrieve--Generate属性转换}
\label{sec:nlp-generation-delete-retrieve}

情感等属性常集中体现在少数短语上：“令人失望”“十分周到”带有明显倾向，
“这家书店的服务”则不带。Delete--Retrieve--Generate（DRG）据此把句子拆成
属性标记与内容两部分：删除原属性的标记，从目标属性语料检索适合当前内容的标记，
再生成新句\scite{li-etal-2018-delete}

\paragraph{属性标记的识别}
设$n_a(u)$为短语$u$在属性$a$语料中的出现次数。DRG用平滑后的相对频率衡量$u$
对属性$a$的\term{属性标记}（attribute marker）程度：
\begin{equation}
  s(u,a)=\frac{n_a(u)+\lambda}{\sum_{a'\neq a}n_{a'}(u)+\lambda},
  \label{eq:nlp-generation-drg-salience}
\end{equation}
$\lambda>0$防止分母为零，并压低只出现一两次的短语。$s(u,a)$超过阈值$\gamma$的
$n$元短语（原文取到四元）被视为属性$a$的标记。
例如“令人失望”在负面语料中出现40次、在正面语料中出现2次，取$\lambda=1$得
$s=41/3\approx13.7$；“服务”在两类语料中分别出现30次和32次，$s=31/33\approx0.94$。
取$\gamma=5$时，只有前者被删除。删除句$x$中属性$a$的全部标记后，
得到内容序列$c(x,a)$，被删的标记序列记为$m(x,a)$。

\paragraph{检索与生成}
检索在目标属性$a'$的语料中寻找内容最接近的句子$x'$：
$x'=\argmin_{x''\in X_{a'}}\operatorname{dist}\bigl(c(x,a),c(x'',a')\bigr)$，
距离可取TF-IDF加权的词重合，或句向量之间的欧氏距离。$x'$的标记$m(x',a')$
就是候选的目标属性表达。由此有四种输出方式：
RetrieveOnly直接返回$x'$；TemplateBased把$m(x',a')$填回$x$中被删除的位置；
DeleteOnly用LSTM编码$c(x,a)$，把编码与目标属性的可学习嵌入拼接作为解码器初态，生成新句；
DeleteAndRetrieve再用另一个编码器读取$m(x',a')$，与内容编码一起交给解码器。

生成器没有跨属性的平行目标，只能用重建训练：以句子自身的$c(x,a)$和属性$a$为输入，
最大化原句$x$的概率。DeleteAndRetrieve的训练输入还包括自身的标记$m(x,a)$，
若原样提供，模型只需学会照抄；因此训练时把它替换为同属性语料中相近的另一组标记，
迫使模型学会把检索来的标记接到内容上。运行时把$a$换成目标属性$a'$、
把$m(x,a)$换成检索到的$m(x',a')$。

以负面到正面的情感转换为例，原句“这家书店的服务令人失望”删除“令人失望”后
留下“这家书店的服务”，正面语料中内容最近的句子是“这家商场的服务十分周到”。
RetrieveOnly直接返回它，把书店换成了商场；TemplateBased把“十分周到”填回空位，
得到“这家书店的服务十分周到”；DeleteOnly只凭内容和属性生成，可能输出“这家书店服务很好”；
DeleteAndRetrieve把检索来的“十分周到”接到内容上，并可调整衔接。
这些是为说明四种方式构造的候选，不是实测输出。

DRG的优点是每一步都可检查：被删的标记、检索到的句子和生成依据都能看到。
它的假设是属性可以由局部短语承载；讽刺、隐含评价或分布在整句结构中的风格不满足这一假设。
按频率识别的标记也可能误伤内容：“须”“不得”这类义务词即使在正式语料中更常见，
也不能当作文体标记删除。

\subsubsection{隐表示解纠缠与回译转换}
\label{sec:nlp-generation-latent-style}

DRG把内容与风格的分离放在词面上。另一条路线把分离放到隐表示里：
若编码器能把句子映射为只含内容的向量$\bm z$，解码器再以目标风格为条件生成，
就不必假设风格集中在可删除的短语。

\paragraph{交叉对齐自编码器}
Shen等考虑两种风格的语料$X_1,X_2$\scite{shen2017crossalign}
编码器$E(x,y)$读取句子$x$及其风格$y$，输出内容表示$\bm z$；生成器$G(y,\bm z)$
是以风格$y$和$\bm z$为初态的GRU。自编码重建损失为
\begin{equation}
  L_{\mathrm{rec}}=\sum_{k=1}^{2}\mathbb{E}_{x\sim X_k}
  \bigl[-\log p_G\bigl(x\mid y_k,E(x,y_k)\bigr)\bigr].
  \label{eq:nlp-generation-shen-rec}
\end{equation}
只有重建并不能阻止$\bm z$携带风格。他们先尝试对齐两种风格下$\bm z$的分布，
用一个判别器区分$\bm z$来自哪种语料，再让编码器骗过它；
进一步的\term{交叉对齐}（cross-alignment）改为对齐生成器的隐状态序列：
判别器$D_1$区分两类序列，一类是生成器以$y_1$为条件、在教师强制下重建真实$x_1$时的隐状态
$\bm h_1$，另一类是把$x_2$的内容$E(x_2,y_2)$配上风格$y_1$、自由生成时的隐状态
$\tilde{\bm h}_2$。若生成器转换得好，后者应与真实风格$y_1$的句子不可区分。
自由生成的离散采样不可求导，他们用带温度的softmax输出作为下一步输入，使梯度可以回传。
$D_2$对称地处理另一方向，总目标为
\begin{equation}
  \begin{aligned}
    &\min_{E,G}\max_{D_1,D_2}\;L_{\mathrm{rec}}
    -\rho\,(L_{\mathrm{adv},1}+L_{\mathrm{adv},2}),\\
    &L_{\mathrm{adv},1}=-\mathbb{E}\bigl[\log D_1(\bm h_1)\bigr]
    -\mathbb{E}\bigl[\log\bigl(1-D_1(\tilde{\bm h}_2)\bigr)\bigr].
  \end{aligned}
  \label{eq:nlp-generation-shen-adv}
\end{equation}
这里$\rho>0$平衡重建与对抗两项。判别器最大化自己的分类对数似然，
编码器和生成器则让转换后的隐状态骗过判别器。这一对抗训练与生成对抗网络相同：
一方生成样本，另一方区分真伪，二者交替更新；其训练不稳定和需要调节更新步数的问题，
可参考\BookChapterRef{chap:neural-generative-models}{卷五的神经生成模型章}。

\paragraph{对抗解纠缠}
Fu等把对抗直接放在内容表示上：一个分类器试图从$\bm z$预测风格，
编码器则被训练成让这个分类器无法预测，再用风格专属的解码器或风格嵌入生成
\scite{fu2018styletransfer}
这类方法统称为解纠缠：希望$\bm z$与风格统计独立，从而“换风格不换内容”。
困难在于，对抗训练使某个分类器失败，并不说明$\bm z$已经不含风格；
训练后另取一个分类器，往往仍能从$\bm z$中读出风格。

\paragraph{去噪与回译}
Lample等的观察正是如此：解纠缠后的表示仍可读出属性，转换却照样有效，
说明解纠缠既难做到，也未必必要\scite{lample2019multiattr}
他们借鉴无监督机器翻译，用两项损失训练一个以属性为条件的编码器—解码器。
去噪自编码项把句子加噪（随机删词、局部打乱）后重建原句；
回译项先用当前模型把$x$转成另一属性$\tilde a$下的$\hat x=G(E(x),\tilde a)$，
不对这一步求导，再要求模型从$\hat x$还原$x$：
\begin{equation}
  L=\lambda_{\mathrm{ae}}\,\mathbb{E}_{x}
    \bigl[-\log p\bigl(x\mid E(\operatorname{noise}(x)),a\bigr)\bigr]
   +\lambda_{\mathrm{bt}}\,\mathbb{E}_{x,\tilde a}
    \bigl[-\log p\bigl(x\mid E(\hat x),a\bigr)\bigr].
  \label{eq:nlp-generation-lample-bt}
\end{equation}
回译项把模型自己生成的$(\hat x,x)$当作伪平行对：输入是属性$\tilde a$的句子，
输出是属性$a$的原句。随着模型变好，伪平行对的质量也提高。
Prabhumoye等用另一种回译：先把英语句子译成法语，再用法语到英语翻译模型的编码器
得到表示，假定翻译过程已削弱了作者风格，最后由各风格专属的解码器生成，
并用风格分类器的反馈训练解码器\scite{prabhumoye2018style}

交叉对齐与对抗解纠缠要求一个难以验证的中间状态，即不含风格的$\bm z$；
去噪加回译放弃这一要求，转而依赖迭代产生的伪平行对，
代价是早期的转换错误会写进伪数据，内容可能逐轮漂移。
两条路线的共同限制是内容保持只受重建损失间接约束，需要单独检查实体、数字与否定。

\subsubsection{大语言模型提示改写}
\label{sec:nlp-generation-llm-style}

前面几种方法都要为每种属性准备语料并单独训练。大语言模型在预训练中见过大量不同风格的文本，
可以直接按指令改写。Reif等提出增强零样本提示：提示中给出几个“原句—某种风格的改写”示例，
示例覆盖多种不同风格而不必包含目标风格，再要求模型把输入改写成目标风格
\scite{reif2022recipe}
示例的作用是让模型理解“只改风格、保留内容”这一任务格式，目标风格本身由一句自然语言描述，
因此“更有诗意”“像新闻标题”这类没有标注语料的风格也能请求。
上下文示例的选择与数目对结果的影响，见第~\ref{sec:nlp-inference-icl-shots}小节。

提示只描述了要求，模型是否做到要靠筛选。具体做法是在提示中列出必须保留的内容，
如“保持人群、日期、时刻与义务不变”；再对同一输入采样$k$个候选，
用风格分类器给出目标属性概率$p_{\mathrm{sty}}(a'\mid\tilde x)$，用句向量给出与源句的相似度$\operatorname{sim}(x,\tilde x)$，
剔除改动了保护片段的候选后，在相似度不低于阈值的候选中取$p_{\mathrm{sty}}$最大者，
或按两者的几何平均排序。这一步就是第~\ref{sec:nlp-inference-reranking}小节的候选重排。
提示改写省去了训练，适用的风格范围也最广；代价是每次调用的推断成本高，
过改与意义漂移需要筛选步骤控制，转换程度也难以像控制码那样精确调节。

\subsection{语法纠错}
\label{sec:nlp-generation-correction}

\term{语法纠错}（grammatical error correction, GEC）的依据是原文存在错误，
而不是另一种表达更好。“请与周五到场”中的同音字、“提高报名流程”中的搭配，
以及“She go to meeting”中的动词词形和冠词，都有明确的修订；
“请于周五到场”已经正确，改成“烦请周五莅临”属于文体选择，不能算作纠错收益。
英文流畅性纠错数据集JFLEG允许比最小语法修改更宽的改写，使用时应沿用其任务定义
\scite{nlp-intro-napoles2017}

纠错有四个挑战。错误稀疏，一句中多数词应原样保留，模型很容易学成“什么都不改”；
真实的错句—正句平行数据有限；一个错句常有多种合法修订；模型又容易过改，
把正确的部分也改掉。方法谱系为：混淆集加语言模型打分；统计翻译式纠错；
神经序列生成；编辑标签模型GECToR；大语言模型。合成错误数据贯穿后三者，
用来弥补平行数据的不足。评价时区分漏改、正确修改、过改与意义改变，
修改级指标见第~\ref{sec:nlp-generation-edit-metrics}小节。

\subsubsection{混淆集与语言分数筛选}
\label{sec:nlp-generation-confusion}

若错误集中在同音、形近或拼写混淆，可以先限制每个位置的替换范围。
\term{混淆集}（confusion set）是某个位置允许尝试的候选字词集合，它来自词典、输入法混淆统计、
人工规则或有标注错误。集合必须包含原词，否则系统遇到正确输入也被迫改写。
给完整候选句计分：
\begin{equation}
  s(y;x)=\log p_{\mathrm{LM}}(y)-\lambda\,d(x,y),
  \label{eq:nlp-generation-confusion-score}
\end{equation}
其中语言模型提供上下文分数，$d$为式~\eqref{eq:nlp-generation-edit-distance}的编辑距离，
$\lambda\geq0$惩罚不必要的改动。只有相对原句的得分增量超过开发集确定的阈值，
且规则检查通过，才执行替换。

对“请／与／周五／到场”，候选位置取$\{\text{与},\text{于},\text{予}\}$。
设三条完整句子的教学对数分数分别为$-8,-4,-10$，$\lambda=1$，
则保留原文得$-8$，替换为“于”得$-5$，替换为“予”得$-11$。
阈值为2时，增量3使“于”通过，再由时间介词规则核验，得到“请于周五到场”。

混淆集方法的覆盖面受限于集合本身：词典没有收录真实误字时，搜索不可能找到正确替换；
逐位置替换也无法补冠词、删重复词或调整远处的搭配。领域术语还须保护，
“束搜索”中的“束”虽不如“数”常见，却不能被改成“数搜索”。
突破逐位置替换的第一步是把纠错看作翻译：错句是源语言，正句是目标语言。
Junczys-Dowmunt与Grundkiewicz用第~\ref{sec:nlp-generation-smt}小节的短语统计翻译实现这一点，
在对数线性模型中加入编辑操作计数等稠密与稀疏特征，并直接按纠错评价指标调节特征权重，
得到了当时的强基线\scite{junczys2016pbmt}
短语表从平行纠错语料中学到“go to meeting→goes to the meeting”这类短语级修改，
不再需要人工列出候选；但它仍受短语表覆盖与局部重排的限制。
后文的神经方法沿用“纠错即翻译”的视角，并用合成数据弥补平行语料的不足。

\subsubsection{合成错误数据}
\label{sec:nlp-generation-synthetic-errors}

序列模型需要大量错句—正句对，真实标注语料却只有几十万到几百万句。
干净文本则几乎无限，若能从正句自动生成逼真的错句，就能构造任意规模的
\term{合成错误数据}（synthetic error data）。

\paragraph{规则加噪}
最直接的做法是对正句随机施加编辑。Grundkiewicz等对每个词元按概率执行替换、删除、
插入或与相邻词交换，其中替换优先从拼写检查器给出的混淆集中取词，
使生成的错误更接近真实拼写与用词错误\scite{grundkiewicz2019synthetic}
算法~\ref{alg:nlp-generation-noise}给出一般形式。

\begin{algorithm}[规则加噪生成错句]
  \label{alg:nlp-generation-noise}
  \begin{algorithmic}[1]
    \Require 正句$y_{1:m}$；错误率分布的均值$\mu$与标准差$\sigma$；
      四种操作的概率$(\pi_{\mathrm{sub}},\pi_{\mathrm{del}},\pi_{\mathrm{ins}},\pi_{\mathrm{swap}})$；
      混淆集$C(\cdot)$与词表$V$
    \Ensure 合成错句$x$
    \State 采样本句错误率$\rho\sim\mathcal{N}(\mu,\sigma^2)$，截断到$[0,1]$；$x\gets y$
    \For{$x$中的每个位置$i$}
      \State 以概率$\rho$决定是否在此处加噪；若不加噪，继续下一个位置
      \State 按$\pi$选择操作：替换为$C(x_i)$中的随机词；删除$x_i$；
        在$x_i$后插入$V$中的随机词；或交换$x_i$与$x_{i+1}$
    \EndFor
    \State \Return $(x,y)$作为训练对
  \end{algorithmic}
\end{algorithm}

$\mu$通常取开发集真实错句相对正句的词级错误率，使合成数据的错误密度与真实数据相近。
例如取四种操作概率为$0.7,0.1,0.1,0.1$、错误率$\rho=0.25$，
对“She goes to the meeting”的期望加噪次数为$5\times0.25=1.25$，
一次可能的结果是把goes换成混淆集中的go、再删除the，得到“She go to meeting”。

\paragraph{反向翻译生成错误}
规则噪声的分布与学习者的真实错误并不相同：真实错误集中在冠词、介词、词形和搭配，
随机插入的词则往往一眼可辨。另一种做法借用机器翻译中的回译：
在真实平行语料上反向训练一个“正句→错句”的模型，再用它把大量干净句子译成错句
\scite{xie2018noising}
若用束搜索解码，反向模型倾向于输出最常见的少数错误，甚至原样复制正句；
Xie等在束搜索的每个候选分数上加入随机噪声或排名惩罚，使生成的错误更多样。
Kiyono等在同一模型结构下比较了直接加噪与回译式合成，以及“先在合成数据上预训练、
再在真实数据上微调”与二者混合训练两种用法，结论是用带噪束搜索的回译数据预训练、
再用真实数据微调最为有效\scite{kiyono2019pseudo}

合成数据的风险在于分布偏移：模型在合成数据上学到的错误类型和错误密度若与真实输入不符，
就会漏改真实错误，或在正确句上“找错”。因此合成数据通常只用于预训练，
再用真实平行数据微调；训练数据中保留一定比例的无错句，模型才会学到“不改”也是正常输出。

\subsubsection{条件序列生成纠错}
\label{sec:nlp-generation-seq-correction}

把纠错当作从错句到正句的翻译，就可以直接使用编码器—解码器整句生成正句。
这条路线能一次写出任意修改，包括长短语插入和重排；问题在于正句的绝大多数词元与源句相同，
逐词元交叉熵大部分花在“照抄”上，真正需要学的修改只占很小比例，模型因此偏向不改。
Junczys-Dowmunt等把纠错看作低资源翻译，提出\term{编辑加权最大似然}
（edit-weighted maximum likelihood），提高被修改词元的损失权重\scite{junczys2018approaching}
\begin{equation}
  L(\bm\theta)=-\sum_{t=1}^{T}\lambda_t\log p_{\bm\theta}(y_t\mid x,y_{<t}),\qquad
  \lambda_t=\begin{cases}\Lambda,& y_t\in\mathcal{M},\\
  1,&y_t\notin\mathcal{M},\end{cases}
  \label{eq:nlp-generation-edit-weighted}
\end{equation}
其中$\mathcal{M}$是目标中被修改的词元，即按式~\eqref{eq:nlp-generation-edit-distance}
与源句对齐后不是原样保留的位置，$\Lambda>1$。
它是式~\eqref{eq:nlp-generation-nll}的加权版本：$\Lambda=1$时退化为普通最大似然，
$\Lambda$越大，模型越敢于修改，精确率也随之下降。
另一种做法从结构上处理“照抄”：Zhao等在Transformer上加入第~\ref{sec:nlp-generation-copy-pointer}小节
的复制分布，并用去噪方式在无标注文本上预训练，使模型把复制当作默认动作，
把生成留给真正需要修改的位置\scite{zhao2019improving}
预训练的BART、T5类编码器—解码器本身就在去噪中学会了复制未损坏部分，
微调到纠错时也继承了这一偏置；其去噪目标见第~\ref{sec:nlp-llm-objective-t5-bart}小节。

整句生成的输出要重新与源句对齐，才能知道模型改了哪里。固定输入
“She go to meeting on Friday”，参考修改为go到goes，以及在to后插入the。
生成候选“She goes to the meeting on Monday”与源句对齐后包含三项编辑：
前两项正确，Friday到Monday是过改，并改变了事实。正确编辑2项、系统编辑3项、
参考编辑2项，精确率为$2/3$，召回率为1。整句生成可以写出任何修改，
也就要为输出中每一处非必要的变化负责。

\subsubsection{GECToR编辑标签纠错}
\label{sec:nlp-generation-gector}

整句生成要逐词重写，推断慢，也给过改留下空间。GECToR沿LaserTagger的思路，
把纠错写成对每个源词元的标签预测，并针对纠错扩充了标签集和推断控制
\scite{omelianchuk-etal-2020-gector}

\paragraph{编码器与两个输出头}
输入是在句首加入特殊位置\code{\$START}的源词元序列，句首位置使得在第一个词之前插入成为可能。
预训练编码器（BERT、RoBERTa或XLNet）给出各位置的向量；一个词被切成多个子词时，
取第一个子词的向量$\bm h_i\in\mathbb{R}^{d}$，基础规模模型$d=768$。
编码器上接两个线性层加softmax：检测头输出该词是否有错的二分类概率
$p_{\mathrm{det}}(\cdot\mid i)=\operatorname{softmax}(\bm W_{\mathrm{d}}\bm h_i+\bm b_{\mathrm{d}})$；
标签头输出变换标签的分布
$p_{\mathrm{tag}}(\cdot\mid i)=\operatorname{softmax}(\bm W_{\mathrm{t}}\bm h_i+\bm b_{\mathrm{t}})$，
$\bm W_{\mathrm{t}}\in\mathbb{R}^{|T|\times d}$。默认标签集共5000个：
4971个基本变换，即\code{\$KEEP}、\code{\$DELETE}和带具体词的\code{\$APPEND\_w}、\code{\$REPLACE\_w}；
另有29个与具体词无关的\term{g变换}（g-transformation），如大小写变换、相邻词合并与拆分、
名词单复数和动词形式变换（\code{\$VERB\_FORM\_VB\_VBZ}把go变成goes）。
g变换让一个标签覆盖一整类词形修改，不必为每个词各设一个替换标签。

训练标签由Levenshtein对齐的一个变体从平行句对导出；一个源词元对应多个修改时，
每轮只保留第一个非\code{\$KEEP}标签，其余修改留给下一轮。检测头的标签由此确定：
标签不是\code{\$KEEP}的位置为“错”。两个头的交叉熵相加作为损失：
\begin{equation}
  L=-\sum_{i=0}^{n}\Bigl[\log p_{\mathrm{tag}}(t_i\mid i)
  +\log p_{\mathrm{det}}(e_i\mid i)\Bigr],\qquad e_i=\mathbb{I}[t_i\neq\code{\$KEEP}].
  \label{eq:nlp-generation-gector-loss}
\end{equation}
检测头与标签头共享编码器，检测头给出的是“这里是否需要动”的整体判断，
推断时可以用来决定整句是否值得修改。

\begin{figure*}[htbp]
  \centering
  % GECToR一轮推断：源词元（含句首位置）经编码器得到逐词向量，检测头与标签头逐位置输出，
% 加置信度偏置后取标签并同时应用，结果进入下一轮。概率为教学示例，与正文算例一致。
\begin{tikzpicture}[
  x=1mm,y=1mm,figure picture,
  tok/.style={figure node,minimum width=17mm,minimum height=6mm,
    inner sep=1pt,draw=FigureInput,fill=FigureInputFill,font=\figurefont\scriptsize},
  enc/.style={figure node,minimum height=7mm,inner sep=1.5pt,
    draw=FigureModel,fill=FigureModelFill,font=\figurefont\scriptsize},
  vec/.style={font=\figurefont\scriptsize},
  outc/.style={figure node,minimum width=19.5mm,minimum height=10mm,align=center,
    inner sep=1pt,draw=FigureOutput,fill=FigureOutputFill,font=\figurefont\scriptsize},
  res/.style={figure node,minimum height=6.5mm,inner sep=1.5pt,
    draw=FigureOutput,fill=FigureOutputFill,font=\figurefont\scriptsize},
  note/.style={font=\figurefont\scriptsize,text=FigureMuted,align=center},
  row label/.style={font=\figurefont\scriptsize,text=FigureMuted,anchor=east,align=right},
  flow/.style={figure flow}
]
  \foreach \i/\w/\tag/\pe in {0/{\$START}/KEEP/0.02,1/She/KEEP/0.03,2/go/{VB$\to$VBZ}/0.85,
      3/to/{APPEND\_the}/0.70,4/meeting/KEEP/0.45}{
    \node[tok] (x\i) at (\i*22,0) {\w};
    \node[vec] (h\i) at (\i*22,22) {$\bm h_{\i}$};
    \node[outc] (t\i) at (\i*22,36) {标签：\tag\\有错概率\pe};
    \draw[flow] (x\i) -- (x\i |- 0,8.5);
    \draw[flow] (\i*22,15.5) -- (h\i);
    \draw[flow] (h\i) -- (t\i);
  }
  \node[enc,minimum width=108mm] (enc) at (44,12)
    {预训练编码器（每个词取首个子词的向量）};
  \node[row label] at (-12,0) {本轮输入};
  \node[row label] at (-12,22) {逐词向量};
  \node[row label] at (-12,36) {标签头\\检测头};

  \node[res,minimum width=60mm] (y) at (44,54) {She goes to the meeting};
  \draw[flow] (t2.north) -- (t2.north |- y.south);
  \draw[flow] (t3.north) -- (t3.north |- y.south);
  \node[note,anchor=west] at (76,54) {同时应用，\\作为下一轮输入};
  \node[row label] at (-12,54) {重建结果};
  \node[note] at (88,44.5) {加偏置$\beta$后KEEP胜出};
\end{tikzpicture}

  \caption{GECToR的一轮推断。每个源词元（含句首位置）取编码器向量，检测头给出有错概率，
  标签头给出变换标签；标签按本轮的原坐标同时应用，得到的新句作为下一轮输入，
  直到不再产生非KEEP标签或达到轮数上限。图中概率为教学示例：meeting处的
  复数变换概率略高于KEEP，加入置信度偏置后保持原词。}
  \label{fig:nlp-generation-gector}
\end{figure*}

\paragraph{推断控制与迭代}
标签头在每个位置给出分布后，GECToR用两个超参数控制修改的积极程度。
\term{置信度偏置}（confidence bias）$\beta\geq0$加到\code{\$KEEP}的概率上，
使模型只在足够确信时才修改；\term{最小错误概率}（minimum error probability）$\delta$
作用于整句：若检测头在所有位置给出的最大有错概率$\max_i p_{\mathrm{det}}(\text{错}\mid i)$
低于$\delta$，整句不改。两者都在开发集上调节，用以权衡精确率与召回率。
如图~\ref{fig:nlp-generation-gector}，设go处
$p(\code{\$VERB\_FORM\_VB\_VBZ})=0.70$、$p(\code{\$KEEP})=0.20$，
meeting处$p(\code{\$NOUN\_NUMBER\_PLURAL})=0.48$、$p(\code{\$KEEP})=0.42$。
取$\beta=0.10$后，go处\code{\$KEEP}升到$0.30$，仍执行词形变换；
meeting处\code{\$KEEP}升到$0.52$，超过复数变换，避免了“meetings”这一过改。

修改之间可能相互依赖，一轮未必能完成。例如“She go to meeting”的目标若是
“She goes to the annual meeting”，每个位置一轮只能追加一个词，
第一轮在to之后追加the，第二轮在新序列的the之后追加annual。
算法~\ref{alg:nlp-generation-gector}给出迭代过程。

\begin{algorithm}[GECToR的迭代纠错]
  \label{alg:nlp-generation-gector}
  \begin{algorithmic}[1]
    \Require 源句$x$；置信度偏置$\beta$，最小错误概率$\delta$，最大轮数$K$
    \Ensure 修订句
    \For{$k=1,\ldots,K$}
      \State 在$x$前加\code{\$START}，编码得到$\bm h_0,\ldots,\bm h_n$
      \State 计算各位置的$p_{\mathrm{det}}$与$p_{\mathrm{tag}}$；
        令$p_{\mathrm{tag}}(\code{\$KEEP}\mid i)\gets p_{\mathrm{tag}}(\code{\$KEEP}\mid i)+\beta$
      \If{$\max_i p_{\mathrm{det}}(\text{错}\mid i)<\delta$}
        \State \Return $x$
      \EndIf
      \State $t_i\gets\argmax_t p_{\mathrm{tag}}(t\mid i)$，$i=0,\ldots,n$
      \If{所有$t_i$都是\code{\$KEEP}}
        \State \Return $x$
      \EndIf
      \State 以本轮$x$的坐标同时应用全部$t_i$，得到新的$x$\Comment{下一轮重新编号}
    \EndFor
    \State \Return $x$
  \end{algorithmic}
\end{algorithm}

一轮中的所有标签都以同一份原句为坐标，从左到右重建，
不能一边插入一边仍按已移动的位置执行后面的动作。原文报告多数修改在前两轮完成，
限制轮数能以少量精度换取速度。

\paragraph{三阶段训练}
GECToR按数据质量分三阶段训练：第一阶段在约九百万句的合成错误数据上预训练；
第二阶段只用真实标注语料中含错误的句子微调；第三阶段在规模较小、质量较高的语料上
用含错句和无错句的混合数据微调。第二阶段集中学习修改，第三阶段把无错句加回来，
让模型学会对正确句不动手，这与上一小节保留无错句的理由相同。
前两个阶段的最初两轮还冻结编码器，只训练两个输出头，避免随机初始化的头把编码器带偏。

\paragraph{与整句生成的比较}
与条件序列生成相比，GECToR的每个位置独立预测，一轮只需一次编码器前向；
原文在批量推断下测得其速度最高可达束宽12的Transformer纠错模型\cite{zhao2019improving}的十倍。
编辑可见，可以逐项接受或拒绝。代价同LaserTagger：标签集决定了可表达的修改，
未收录的替换、长短语插入和大幅重排需要多轮，甚至无法完成。
它的g变换是为英语形态设计的；用于中文时，需要重新规定按字还是按词对齐、
同音字替换动作和词语边界，而不能把单复数、时态标签原样搬过去。

\subsubsection{大语言模型纠错}
\label{sec:nlp-generation-llm-gec}

大语言模型可以直接按指令纠错：给出错句和“修正其中的语法错误”的要求，
或在纠错平行语料上微调。它的优势是无需专门的标签集，能处理搭配、语序与上下文相关的错误，
输出也更流畅。Fang等对ChatGPT的系统评测发现，它在允许流畅改写的JFLEG上表现突出，
在以最小编辑为参考的数据集上自动分数却明显低于专门的纠错系统；
人工分析显示其主要问题是\term{过度修改}（overcorrection）：
把语法正确但不够地道的部分也改写，与最小编辑参考对不上\scite{fang2023chatgpt}

过度修改既是评价口径问题，也是实际风险：改写可能改动实体、数字或语气。
控制方法有两类：一是在提示中明确要求“只修正语法错误，尽量少改，不改动专名和数字”，
并给出最小编辑的示例；二是事后过滤：用式~\eqref{eq:nlp-generation-edit-distance}
把输出与源句对齐，恢复编辑集合，删去落在保护片段上的编辑和不属于目标错误类型的编辑，
再把保留的编辑应用到源句。后一类把大语言模型当作候选编辑的提供者，
最终修改仍可逐项检查。评价时同时报告最小编辑与流畅性两类参考上的结果，
具体指标见第~\ref{sec:nlp-generation-edit-metrics}小节。

\subsection{文本简化}
\label{sec:nlp-generation-simplification}

\term{文本简化}（text simplification）要为指定读者降低理解负担，
先确定必须保留的信息和可删的背景，再通过词汇替换、拆句、删除和重排实现。
它的挑战是：平行简化语料规模有限，常见的维基百科与简明维基百科句对对齐还有噪声；
“更短”“用更常见的词”只是代理属性，删得越多越短，却可能删掉关键条件；
不同读者需要的简化程度不同，一个模型只学到训练语料的平均程度。
ASSET的多参考简化同时包含词汇替换、拆句和删减，提醒评价不能只奖励缩短
\scite{nlp-intro-alvamanchego2020}
方法谱系为：规则与统计翻译式简化；以简化指标为奖励的强化学习；
用控制属性显式调节简化程度；从无标注语料挖掘复述实现无监督简化。
自动指标SARI见第~\ref{sec:nlp-generation-sari}小节，事实保持与读者理解则需要另行检查。

\subsubsection{词典替换、拆句与规则重排}
\label{sec:nlp-generation-rule-simplification}

规则路线从适合读者的词典中取常用表达替换难词，再利用句法结构确定拆句位置和省略成分。
例如把“须”解释为“必须”，而不替换成义务更弱的“可以”；拆开条件句时，把适用人群带到新句中；
代词只有在指代唯一时才换成明确名词。句法与指代分析见第~\ref{chap:nlp-text-analysis}章。

以带否定与条件的长句为例：
“若您是校外来宾，须在4月16日18:00前提交报名信息；这项要求不适用于校内师生。”
分析得到第一分句的条件是“校外”，主干是“须提交”；第二分句的“这项要求”指提前提交报名信息。
替换“须”、把条件改为显式主语并拆句，得到
“校外来宾必须在4月16日18:00前提交报名信息。校内师生不受这项要求限制。”
若改成“请在4月16日18:00前提交报名信息。校内师生不受限制”，
第一句丢掉了人群范围，第二句把一项要求扩大成所有限制。拆句减少了句内依赖，
却增加了跨句指代；保留原指代比补错指代更安全。

规则可审查，但覆盖有限。统计翻译式简化把复杂句到简单句看作单语翻译：
Wubben等在维基百科—简明维基百科句对上训练短语统计翻译模型PBMT-R
\scite{wubben2012sentence}
简化语料中大量句对几乎相同，模型学到的最优输出常常就是原句；
他们因此在模型分数最高的$N$个候选$\mathcal{N}_N(x)$中，
选与原句差异最大的那一个：
\begin{equation}
  \hat y=\argmax_{y\in\mathcal{N}_N(x)}d(x,y),
  \label{eq:nlp-generation-pbmtr}
\end{equation}
$d$为词级编辑距离。这一重排强迫模型做出修改，但“与原句差异大”并不等于“更简单”，
可能选中删错或改错的候选。下面的DRESS把“更简单”本身写进训练目标。

\subsubsection{DRESS强化学习简化}
\label{sec:nlp-generation-dress}

最大似然训练让模型模仿参考简化句，而参考中大量是原样复制，
模型学到的往往是保守的复制。Zhang与Lapata提出DRESS，
直接以简化质量为奖励训练生成模型\scite{nlp-generation-zhang2017dress}
策略是一个带注意力的LSTM编码器—解码器$p_{\bm\theta}(y\mid x)$，
先用最大似然预训练，再从中采样完整输出$\hat y$，按三项奖励的加权和打分：
\begin{equation}
  r(\hat y)=\lambda_{\mathrm{S}}\operatorname{SARI}(x,\hat y,y^{\ast})
  +\lambda_{\mathrm{R}}\cos(\bm q_x,\bm q_{\hat y})
  +\lambda_{\mathrm{F}}\exp\Bigl(\frac{1}{|\hat y|}\sum_{t}\log p_{\mathrm{LM}}(\hat y_t\mid\hat y_{<t})\Bigr).
  \label{eq:nlp-generation-dress-reward}
\end{equation}
第一项是简化指标SARI，它同时比较源句、输出和参考，奖励正确的增加、保留与删除；
第二项是意义相关性，$\bm q_x,\bm q_{\hat y}$为一个预训练序列自编码器给出的句向量；
第三项是流畅度，即简单语料上训练的语言模型给出的逐词元概率的几何平均，
也就是困惑度的倒数。三项分别对应简单、保意和通顺，单独优化任一项都会退化：
只看SARI可能删掉关键信息，只看相关性会原样复制，只看流畅度会生成短而空洞的句子。

参数用REINFORCE加基线更新，梯度为
\[
  \nabla_{\bm\theta}J\approx\bigl(r(\hat y)-b(x)\bigr)\nabla_{\bm\theta}\log p_{\bm\theta}(\hat y\mid x),
\]
基线$b(x)$由解码器隐状态的线性回归估计，用于降低梯度方差。
这正是第~\ref{sec:nlp-generation-exposure-bias}小节序列级训练的一个实例：
训练时模型以自己的采样为前缀，奖励直接对应任务目标。
设教学权重$\lambda_{\mathrm{S}}=1$、$\lambda_{\mathrm{R}}=\lambda_{\mathrm{F}}=0.5$，
某个采样的SARI为0.40、相关性0.80、流畅度0.50，则$r=0.40+0.40+0.25=1.05$；
基线为0.95时优势为0.10，这条采样的对数概率被小幅提高。
另一条采样删掉了截止时间，SARI升到0.45，相关性却降到0.50，
$r=0.45+0.25+0.25=0.95$，优势为零，不受鼓励。

DRESS直接对准简化目标，但三项权重需要在开发集上调节，奖励函数本身也只是代理：
SARI依赖参考，相关性与流畅度来自各自的模型，强化学习可能找到骗过它们的输出。

\subsubsection{ACCESS受控简化}
\label{sec:nlp-generation-access}

DRESS学到的简化程度是训练数据与奖励权重决定的平均值，用户无法指定“保留信息但多拆句”。
ACCESS把目标属性编码为输入前的控制标记，让同一个模型按指定程度简化
\scite{martin-etal-2020-controllable}
它使用四类可由句对直接计算的属性：字符长度比、字符级编辑相似度、
词汇复杂度和依存树深度。训练时对每个句对测出这四个值，离散化后作为前缀
拼到原句前，模型学习“在这些属性下，原句应怎样变成简化句”。

记原句与目标的字符数为$n_x,n_y$，长度比为$n_y/n_x$；编辑相似度为
$1-d_{\mathrm{char}}(x,y)/\max(n_x,n_y)$，值越大表示改动越少；
词汇复杂度先取各词词频排名的对数，用第三四分位数汇总，再取目标与源之比；
句法复杂度取依存分析得到的最大依存深度，本节统一为目标深度与源深度之比。
例如一对教学句经固定预处理后有40与36个字符、字符编辑距离12，
对数词频排名的第三四分位数为5与4，最大依存深度为5与3，
四个控制值便是$0.90,0.70,0.80,0.60$。按0.05为间隔离散并把比值上限设为2，
把对应的四个控制标记接在原句前，以真实简化句为目标训练条件最大似然。

推断时目标句未知，控制值由用户或开发集选定，例如要求长度比0.8、编辑相似度0.6，
模型便倾向于输出更短、改动更多的句子。控制标记是条件而不是硬约束：
生成后应重新计算实际属性，检查是否接近请求值；属性达标也不说明截止时间、
义务和人群都已保留，这些仍要逐项核对。

\subsubsection{MUSS挖掘复述的无监督简化}
\label{sec:nlp-generation-muss}

ACCESS仍需平行简化语料，而这类语料在英语之外很少。MUSS的观察是：
在大量复述对中，有的方向是变简单，有的是变复杂，长度和词频都在变化；
若用ACCESS的控制标记把每对的属性变化写进输入，模型就能从普通复述中学会按控制值改写，
推断时把控制值设到“更简单”一侧即可\scite{martin2022muss}
这样，简化所需的监督换成了更容易获得的复述。

复述对从网络语料中挖掘：把大规模单语网页文本切成句子或短句序列，
用多语句向量模型LASER编码，以近似最近邻检索为每个序列找最相近的若干序列，
保留距离低于阈值的对；再过滤掉几乎相同的对和一方包含另一方的对，
这些对没有改写信息。挖掘结果按ACCESS的方式计算控制值，在其上微调预训练的BART
（多语时用mBART）。推断时控制值在开发集上选择。
这一做法把第~\ref{sec:nlp-generation-backtranslation-paraphrase}小节的复述构造从翻译换成检索，
也是第~\ref{chap:nlp-classification-matching-organization}章向量检索在数据构造中的用法。

MUSS不需要简化标注，可用于标注稀缺的语言；代价是挖掘与过滤的计算量、
仍需开发集或人工设定的控制值，以及挖掘到的复述本身可能改变意义，这类错误会传给模型。

\paragraph{文本编辑方法的对照}
表~\ref{tab:nlp-generation-editing-methods}按四个任务汇总本节方法。
纵向看，输出表示从局部替换、整句生成到编辑标签，再到提示，
越靠近整句生成，可表达的修改越多，过改的风险也越大；
编辑标签和选择式方法则以有限的表达能力换取速度和可检查性。
横向看，平行数据不足是共同问题，回译、合成错误、属性语料和挖掘复述
是四种替代监督，各自带来分布偏移或意义漂移，需要过滤与核验配合。

\begin{table*}[htbp]
  \centering
  \small
  \begin{tabularx}{\linewidth}{@{}>{\raggedright\arraybackslash}p{24mm}
    >{\raggedright\arraybackslash}p{30mm}>{\raggedright\arraybackslash}p{38mm}
    >{\raggedright\arraybackslash}p{30mm}>{\raggedright\arraybackslash}X@{}}
    \toprule
    方法 & 解决的问题 & 核心机制 & 数据需求 & 新代价与风险\\
    \midrule
    回译复述 & 复述对稀缺 & 经枢轴语言译回，按相似度过滤 & 翻译系统或双语语料 & 意义漂移只能事后过滤\\
    LaserTagger & 整句重写慢、易补造 & KEEP/DELETE加词表内插入短语 & 平行句对 & 词表外插入与大幅重排无法表达\\
    DRG & 无平行风格数据 & 删属性标记、检索、重建 & 按属性分开的语料 & 依赖属性可由局部短语承载\\
    交叉对齐、对抗解纠缠 & 风格不局限于短语 & 对抗对齐隐表示或隐状态 & 属性语料 & 解纠缠难以验证，训练不稳\\
    去噪加回译 & 不必解纠缠 & 迭代生成伪平行对 & 属性语料 & 伪数据错误逐轮累积\\
    提示改写 & 每种风格都要训练 & 示例加风格描述，候选筛选 & 无需训练 & 推断成本高，过改需筛选\\
    混淆集打分 & 同音形近错误 & 局部候选加语言分数减编辑代价 & 混淆集与语言模型 & 覆盖受限于候选集合\\
    合成错误数据 & 平行纠错数据不足 & 规则加噪或反向模型生成错句 & 干净文本 & 合成错误分布偏移\\
    序列生成纠错 & 任意形式的修改 & 编码器—解码器，编辑加权或复制 & 平行数据加合成数据 & 推断慢，需对齐恢复编辑\\
    GECToR & 推断慢，过改 & 检测头加标签头，偏置与阈值，迭代 & 三阶段数据 & 标签集限制表达，英语专用变换\\
    大模型纠错 & 专门标签与训练成本 & 提示或微调，事后过滤编辑 & 少量或无 & 过度修改\\
    PBMT-R & 规则覆盖有限 & 单语短语翻译，按差异重排 & 平行简化语料 & 差异大不等于更简单\\
    DRESS & 最大似然偏向复制 & 以SARI、相关性、流畅度为奖励 & 平行简化语料 & 奖励权重需调，代理可被利用\\
    ACCESS & 简化程度不可控 & 四类属性的离散控制标记 & 平行简化语料 & 控制值是条件而非约束\\
    MUSS & 简化语料稀缺 & 挖掘复述加控制标记 & 无标注单语语料 & 挖掘成本，复述本身可能漂移\\
    \bottomrule
  \end{tabularx}
  \caption{文本编辑方法对照。“数据需求”指训练该方法所需的主要监督来源；
  各方法的输出都需要按任务另行核验保护片段与意义是否保持。}
  \label{tab:nlp-generation-editing-methods}
\end{table*}
\FloatBarrier
